% Part: sets-functions-relations
% Chapter: infinite
% Section: dedekind-induction

\documentclass[../../../include/open-logic-section]{subfiles}

\begin{document}

\olfileid{sfr}{infinite}{induction}
\olsection[算术归纳法]{Dedekind 代数与算术归纳法}

现在关键的一点是，一个 Dedekind 代数——事实上，\emph{任意} Dedekind 代数——都将充当自然数的替代品。这得益于以下平凡推论：

\begin{thm}[算术归纳法]\ollabel{thm:dedinfiniteinduction}
设 $N, s, o$ 构成一个 Dedekind 代数。则对任意集合 $X$：  
	\begin{center}
		若 $o \in X$ 且 $(\forall {n} \in N \cap X){s}({n}) \in X$，{则} $N \subseteq X$。
	\end{center}
\end{thm}

\begin{proof}
由 Dedekind 代数的定义，$N = \closureofunder{s}{o}$。现在若 ${o} \in X$ 且 $(\forall {n} \in N)(n \in X \lif
{s}({n}) \in X)$，则 $N = \closureofunder{s}{o} \subseteq X$。
\end{proof}

既然归纳法是自然数的特征，要点便在于：给定任意 Dedekind 无穷集，我们都可以构造一个 Dedekind 代数，并将该代数用作自然数的替代品。

诚然，\olref{thm:dedinfiniteinduction} 是以\emph{集合论}术语表述归纳法的。但我们很容易将这一原理用更熟悉的方式陈述出来：

\begin{cor}\ollabel{natinductionschema}
设 $N, s, o$ 构成一个 Dedekind 代数。则对任一可带参数的公式 $\phi(x)$，有：
\begin{center}
	若 $\phi(o)$ 且 $(\forall {n} \in N)(\phi(n)\lif
	\phi({s}({n})))$，{则} $(\forall n \in N)\phi(n)$
\end{center}
\end{cor}

\begin{proof}
令 $X = \Setabs{n \in N}{\phi(n)}$，然后应用 \olref{thm:dedinfiniteinduction}。
\end{proof}

在这一结果中，我们提到公式“带有参数”。粗略地说，这意味着：对任意对象 $c_1, \ldots, c_k$，我们可以处理 $\phi(x, c_1, \ldots, c_k)$。更精确地说，我们可以无需提及“参数”来陈述这一结果如下：对任意公式 $\phi(x, v_1, \ldots, v_k)$（其自由变元均已列出），我们有：
	\begin{align*}
			\forall v_1 \ldots \forall v_k((&\phi(o, v_1,\ldots, v_k) \land {}\\
			&	(\forall x \in N)(\phi(x,v_1, \ldots, v_k) \lif \phi(s(x), v_1,\ldots, v_k))) \lif {}\\
			&\hspace{3em} (\forall x \in N)\phi(x, v_1,\ldots, v_k))
	\end{align*}
显然，使用“带有参数”的说法可以使陈述更易读。（在 \olref[sth][][]{part} 中，我们将相当频繁地使用这一方法。）

回到 Dedekind 代数：给定任意一个 Dedekind 代数，我们也可以定义通常的算术函数，如加法、乘法和乘方。然而这并非显然，它涉及\emph{递归定义}的技术。我们将在更靠后且更一般的语境中引入该技术并论证其合理性。（有兴趣的读者不妨在读完 \olref[sth][ord-arithmetic][]{chap} 后回顾此处，或者阅读另一种处理方式，例如 \citealt[pp.~95--8]{Potter2004}。）但是，当 $N, s, o$ 构成一个 Dedekind 代数时，我们最终将能作出如下规定：
\begin{align*}
	{a} + {o} &= {a} & & & {a} \times {o} &= {o} & & & {a}^{o} &= s(o)\\
{a} + {s}({b}) &= {s}({a}+{b}) &&& {a} \times {s}({b}) &= ({a}\times {b}) + {a}  & & & {a}^{{s}({b})} &= {a}^{b} \times {a}
\end{align*}
并证明其行为符合预期。

\end{document}